% GENERATED FILE. Edit canonical lessons, then run npm run generate:books.
% canonical-source-hash: fnv1a64:c9989c0c4de960fe
% canonical-lessons: TA-C478-uruppu, TA-C478-nuraiyiral, TA-C478-kalliral, TA-C478-cirunirakam, TA-C478-kutal

\chapter{Inside the Body}
\label{ch:ta-inside-the-body}
\hlchaptermodality{478}

\begin{chapteropening}
\textbf{By the end of this chapter:} \emph{I can name organs such as the lungs, the liver and the kidneys.}

Complete the last lesson of chapter 478: I can name organs such as the lungs, the liver and the kidneys.
\end{chapteropening}

\section[\texorpdfstring{\ta{உறுப்பு} (uṟuppu)}{உறுப்பு (uṟuppu)}]{\ta{உறுப்பு} (uṟuppu) — an organ; a limb}

\label{lesson:TA-C478-uruppu}



\begin{quote}
Before the new one: say the Tamil for asthma, then the Tamil for jaundice.
\end{quote}

\subsection*{You'll want to know: \ta{உறுப்பு}}
\textbf{\ta{உறுப்பு}} — \emph{uṟuppu} — \textquotedblleft{}an organ; a limb\textquotedblright{}.

\begin{quote}
\ta{இதயம்} \ta{உடலின்} \ta{முக்கியமான} \ta{உறுப்பு}.
\end{quote}

\begin{quote}
\emph{itayam uṭaliṉ mukkiyamāṉa uṟuppu.}
\end{quote}

\begin{quote}
The heart is an important organ of the body.
\end{quote}

Plural: \ta{உறுப்புகள்}, \emph{uṟuppukaḷ}.

Also a member of a group: \ta{குழு} \ta{உறுப்பினர்}.

\begin{cousinweb}
\textbf{\ta{உறுப்பு}} (\emph{uṟuppu}) is native Dravidian: DEDR 710 lists it, \textquotedblleft{}body, limb\textquotedblright{}, in the family of \ta{உறு} \emph{uṟu}, \textquotedblleft{}to be joined\textquotedblright{}.
\end{cousinweb}

\begin{grammarlens}[title={What you've built}]
One more word for describing what you lived and felt.
\end{grammarlens}

\subsection*{Your turn}
\begin{itemize}
\raggedright
  \item \emph{Say it:} \emph{uṟuppu}
  \item \emph{Say it:} \emph{uṟuppu}, once more
  \item \emph{Say it:} \emph{itayam uṭaliṉ mukkiyamāṉa uṟuppu.} — The heart is an important organ of the body.
  \item \emph{Recall:} say the Tamil for asthma, then the Tamil for jaundice, then say \emph{uṟuppu} again
\end{itemize}

\begin{tcolorbox}[breakable,colback=teal!4,colframe=teal!35!black,title={Before you move on}]
What does \ta{உறுப்பு} mean? (\textquotedblleft{}An organ; a limb\textquotedblright{}.) Say it once more.
\end{tcolorbox}
\section[\texorpdfstring{\ta{நுரையீரல்} (nuraiyīral)}{நுரையீரல் (nuraiyīral)}]{\ta{நுரையீரல்} (nuraiyīral) — a lung}

\label{lesson:TA-C478-nuraiyiral}



\begin{quote}
Before the new one: say the Tamil for jaundice, then the Tamil for an organ.
\end{quote}

\subsection*{You'll want to know: \ta{நுரையீரல்}}
\textbf{\ta{நுரையீரல்}} — \emph{nuraiyīral} — \textquotedblleft{}a lung\textquotedblright{}.

\begin{quote}
\ta{புகை} \ta{நுரையீரலுக்குக்} \ta{கெடுதல்}.
\end{quote}

\begin{quote}
\emph{pukai nuraiyīralukkuk keṭutal.}
\end{quote}

\begin{quote}
Smoke is harmful to the lungs.
\end{quote}

Literally the frothy inner organ.

\begin{cousinweb}
\textbf{\ta{நுரையீரல்}} (\emph{nuraiyīral}) is a native compound of \ta{நுரை} \emph{nurai}, \textquotedblleft{}froth\textquotedblright{} (DEDR 3710), and \emph{īral}, \textquotedblleft{}an inner organ, the liver\textquotedblright{} (DEDR 546).
\end{cousinweb}

\begin{grammarlens}[title={What you've built}]
One more word for describing what you lived and felt.
\end{grammarlens}

\subsection*{Your turn}
\begin{itemize}
\raggedright
  \item \emph{Say it:} \emph{nuraiyīral}
  \item \emph{Say it:} \emph{nuraiyīral}, once more
  \item \emph{Say it:} \emph{pukai nuraiyīralukkuk keṭutal.} — Smoke is harmful to the lungs.
  \item \emph{Recall:} say the Tamil for jaundice, then the Tamil for an organ, then say \emph{nuraiyīral} again
\end{itemize}

\begin{tcolorbox}[breakable,colback=teal!4,colframe=teal!35!black,title={Before you move on}]
What does \ta{நுரையீரல்} mean? (\textquotedblleft{}A lung\textquotedblright{}.) Say it once more.
\end{tcolorbox}
\section[\texorpdfstring{\ta{கல்லீரல்} (kallīral)}{கல்லீரல் (kallīral)}]{\ta{கல்லீரல்} (kallīral) — the liver}

\label{lesson:TA-C478-kalliral}



\begin{quote}
Before the new one: say the Tamil for an organ, then the Tamil for a lung.
\end{quote}

\subsection*{You'll want to know: \ta{கல்லீரல்}}
\textbf{\ta{கல்லீரல்}} — \emph{kallīral} — \textquotedblleft{}the liver\textquotedblright{}.

\begin{quote}
\ta{அதிகமாகக்} \ta{குடிப்பது} \ta{கல்லீரலைக்} \ta{கெடுக்கும்}.
\end{quote}

\begin{quote}
\emph{atikamākak kuṭippatu kallīralaik keṭukkum.}
\end{quote}

\begin{quote}
Drinking too much harms the liver.
\end{quote}

The doubled \ta{ல்ல்} of \ta{கல்} + a vowel.

\begin{cousinweb}
\textbf{\ta{கல்லீரல்}} (\emph{kallīral}) is native Dravidian: its second part is \emph{īral}, \textquotedblleft{}an inner organ, the liver\textquotedblright{} (DEDR 546); the first part looks like \ta{கல்} \textquotedblleft{}stone\textquotedblright{}, but that reading is not settled.
\end{cousinweb}

\begin{grammarlens}[title={What you've built}]
One more word for describing what you lived and felt.
\end{grammarlens}

\subsection*{Your turn}
\begin{itemize}
\raggedright
  \item \emph{Say it:} \emph{kallīral}
  \item \emph{Say it:} \emph{kallīral}, once more
  \item \emph{Say it:} \emph{atikamākak kuṭippatu kallīralaik keṭukkum.} — Drinking too much harms the liver.
  \item \emph{Recall:} say the Tamil for an organ, then the Tamil for a lung, then say \emph{kallīral} again
\end{itemize}

\begin{tcolorbox}[breakable,colback=teal!4,colframe=teal!35!black,title={Before you move on}]
What does \ta{கல்லீரல்} mean? (\textquotedblleft{}The liver\textquotedblright{}.) Say it once more.
\end{tcolorbox}
\section[\texorpdfstring{\ta{சிறுநீரகம்} (ciṟunīrakam)}{சிறுநீரகம் (ciṟunīrakam)}]{\ta{சிறுநீரகம்} (ciṟunīrakam) — a kidney}

\label{lesson:TA-C478-cirunirakam}



\begin{quote}
Before the new one: say the Tamil for a lung, then the Tamil for the liver.
\end{quote}

\subsection*{You'll want to know: \ta{சிறுநீரகம்}}
\textbf{\ta{சிறுநீரகம்}} — \emph{ciṟunīrakam} — \textquotedblleft{}a kidney\textquotedblright{}.

\begin{quote}
\ta{நிறைய} \ta{தண்ணீர்} \ta{குடிப்பது} \ta{சிறுநீரகத்துக்கு} \ta{நல்லது}.
\end{quote}

\begin{quote}
\emph{niṟaiya taṇṇīr kuṭippatu ciṟunīrakattukku nallatu.}
\end{quote}

\begin{quote}
Drinking plenty of water is good for the kidneys.
\end{quote}

\ta{சிறுநீர்} \textquotedblleft{}urine\textquotedblright{} + \ta{அகம்} \textquotedblleft{}place\textquotedblright{}.

\begin{cousinweb}
\textbf{\ta{சிறுநீரகம்}} (\emph{ciṟunīrakam}) is a native compound built from \ta{சிறு} \emph{ciṟu}, \textquotedblleft{}small\textquotedblright{} (DEDR 1594), \ta{நீர்} \emph{nīr}, \textquotedblleft{}water\textquotedblright{} (DEDR 3690), and \ta{அகம்}, \textquotedblleft{}place\textquotedblright{}.
\end{cousinweb}

\begin{grammarlens}[title={What you've built}]
One more word for describing what you lived and felt.
\end{grammarlens}

\subsection*{Your turn}
\begin{itemize}
\raggedright
  \item \emph{Say it:} \emph{ciṟunīrakam}
  \item \emph{Say it:} \emph{ciṟunīrakam}, once more
  \item \emph{Say it:} \emph{niṟaiya taṇṇīr kuṭippatu ciṟunīrakattukku nallatu.} — Drinking plenty of water is good for the kidneys.
  \item \emph{Recall:} say the Tamil for a lung, then the Tamil for the liver, then say \emph{ciṟunīrakam} again
\end{itemize}

\begin{tcolorbox}[breakable,colback=teal!4,colframe=teal!35!black,title={Before you move on}]
What does \ta{சிறுநீரகம்} mean? (\textquotedblleft{}A kidney\textquotedblright{}.) Say it once more.
\end{tcolorbox}
\section[\texorpdfstring{\ta{குடல்} (kuṭal)}{குடல் (kuṭal)}]{\ta{குடல்} (kuṭal) — the intestines; the gut}

\label{lesson:TA-C478-kutal}



\begin{quote}
Before the new one: say the Tamil for the liver, then the Tamil for a kidney.
\end{quote}

\subsection*{You'll want to know: \ta{குடல்}}
\textbf{\ta{குடல்}} — \emph{kuṭal} — \textquotedblleft{}the intestines; the gut\textquotedblright{}.

\begin{quote}
\ta{காரமான} \ta{உணவு} \ta{சிலருக்குக்} \ta{குடலைக்} \ta{கெடுக்கும்}.
\end{quote}

\begin{quote}
\emph{kāramāṉa uṇavu cilarukkuk kuṭalaik keṭukkum.}
\end{quote}

\begin{quote}
Spicy food upsets the gut for some people.
\end{quote}

\ta{சிறுகுடல்}, the small intestine.

\begin{cousinweb}
\textbf{\ta{குடல்}} (\emph{kuṭal}) is native Dravidian (DEDR 1652); Malayalam has the same word.
\end{cousinweb}

\begin{grammarlens}[title={What you've built}]
One more word for describing what you lived and felt.
\end{grammarlens}

\subsection*{Your turn}
\begin{itemize}
\raggedright
  \item \emph{Say it:} \emph{kuṭal}
  \item \emph{Say it:} \emph{kuṭal}, once more
  \item \emph{Say it:} \emph{kāramāṉa uṇavu cilarukkuk kuṭalaik keṭukkum.} — Spicy food upsets the gut for some people.
  \item \emph{Recall:} say the Tamil for the liver, then the Tamil for a kidney, then say \emph{kuṭal} again
\end{itemize}

\begin{tcolorbox}[breakable,colback=teal!4,colframe=teal!35!black,title={Before you move on}]
What does \ta{குடல்} mean? (\textquotedblleft{}The intestines; the gut\textquotedblright{}.) Say it once more.
\end{tcolorbox}
